% !TEX root = ../../main.tex

\subsection{Tree Representation and Inferential Objective}

The method begins with a hierarchical clustering tree built from the samples.
The tree should be read as a structured list of candidate groups and candidate
splits, not as a claim that every branch already corresponds to a real cluster.
The formal node, edge, and descendant-set definitions are given in the next
subsections.

Once the tree is understood in that way, the inferential problem becomes
concrete. Each internal node proposes a possible split of the sample set into
its two child subtrees, but the tree alone does not say whether that split
should be kept in the final partition. The main question is whether a proposed
split is supported strongly enough to remain as a cluster boundary. The role of
the statistical method is therefore not to build the hierarchy from scratch,
but to evaluate the splits proposed by that hierarchy.

KL-TE implements the split decision by separating representation, testing, and
traversal. The representation step maps the samples into one working
sample-feature matrix and summarizes every tree node by the empirical feature
distribution of its descendant leaves. The testing step attaches child-parent
and sibling evidence to the candidate split. The traversal step reads those
annotations and returns one terminal partition.

This structure is needed because the tree creates statistical issues that a
purely geometric cut rule does not address. Parent and child are nested groups,
so the edge comparison is not an ordinary two-sample problem. The feature space
can be high-dimensional and strongly correlated, so evidence has to be combined
carefully rather than coordinate by coordinate. Because many candidate splits
are tested across the tree, multiplicity must be controlled. In addition, the
sibling test is affected by post-selection scale distortion, since the same data are
used both to build the hierarchy and to evaluate the splits that the hierarchy
exposes. Here post-selection scale distortion means that a split can look more
statistically convincing than it really is because it was selected from the
same data on which it is later tested. The rest of this section develops the
method in that order.

Accordingly, the next subsection defines the representation and subtree
objects. The following subsections then introduce the edge test, the edge
multiplicity correction, the sibling test, and the traversal rule.

\begin{figure}[p]
\centering
\resizebox{0.77\textwidth}{!}{\begin{tikzpicture}[
  x=1cm,
  y=1cm,
  font=\small,
  >=Latex,
]
\tikzset{
  topanel/.style={
    draw,
    rounded corners=2pt,
    fill=gray!4,
    minimum width=4.60cm,
    minimum height=4.95cm,
    inner sep=5pt,
  },
  widepanel/.style={
    draw,
    rounded corners=2pt,
    fill=gray!4,
    minimum width=9.75cm,
    minimum height=4.75cm,
    inner sep=5pt,
  },
  panel/.style={
    draw,
    rounded corners=2pt,
    fill=gray!4,
    minimum width=6.10cm,
    minimum height=4.20cm,
    inner sep=4pt,
  },
  panelsmall/.style={
    draw,
    rounded corners=2pt,
    fill=gray!4,
    minimum width=4.00cm,
    minimum height=4.75cm,
    inner sep=4pt,
  },
  paneltitle/.style={font=\bfseries\footnotesize, anchor=north west},
  steparrow/.style={->, line width=0.88pt, draw=black!62, shorten <=2pt, shorten >=2pt},
  treeline/.style={line width=0.9pt, draw=black!75},
  helperline/.style={draw=black!28, line width=0.4pt},
  gate2line/.style={draw=orange!98!black, line width=1.45pt},
  gate3line/.style={draw=green!75!black, line width=1.45pt},
  cutline/.style={draw=red!75!black, line width=1.25pt, dashed},
  treept/.style={circle, draw=black!75, fill=white, minimum size=3.9mm, inner sep=0pt},
  tinylabel/.style={font=\scriptsize, inner sep=1pt},
  leafchip/.style={
    draw=black!18,
    rounded corners=1pt,
    fill=white,
    inner sep=1.1pt,
    font=\tiny,
    align=center,
    text width=1.14cm,
  },
  summarychip/.style={
    draw=black!22,
    rounded corners=1pt,
    fill=white,
    inner sep=1.6pt,
    font=\scriptsize,
    align=center,
    text width=1.72cm,
  },
  eqbox/.style={
    draw=black!25,
    rounded corners=1pt,
    fill=white,
    inner sep=2.5pt,
    font=\scriptsize,
    align=center,
    text width=3.65cm,
  },
  gate2box/.style={eqbox, draw=orange!98!black, fill=orange!8},
  gate3box/.style={eqbox, draw=green!75!black, fill=green!8}
}


% Top block: cards 1--3 keep their existing relative arrangement.
\begin{scope}[shift={(0,0)}]
  \node[topanel] (p1) at (0,0) {};
\node[paneltitle] at ($(p1.north west)+(0.10,-0.12)$) {1. Input matrix};
\node[
  anchor=north west,
  font=\scriptsize,
  align=left,
  text width=4.00cm,
] at ($(p1.north west)+(0.10,-0.38)$)
  {discrete sample-feature matrix $X \in \{0,1\}^{n \times p}$};

\matrix[
  matrix of nodes,
  nodes={draw, minimum width=0.50cm, minimum height=0.50cm, anchor=center, font=\scriptsize},
  row sep=-\pgflinewidth,
  column sep=-\pgflinewidth,
] (mat) at ($(p1.center)+(0,0.45)$) {
1 & 1 & 0 & 1 \\
1 & 1 & 0 & 1 \\
0 & 0 & 1 & 0 \\
0 & 0 & 1 & 0 \\
};
\node[tinylabel, anchor=east] at ($(mat-1-1.west)+(-0.12,0)$) {$A$};
\node[tinylabel, anchor=east] at ($(mat-2-1.west)+(-0.12,0)$) {$B$};
\node[tinylabel, anchor=east] at ($(mat-3-1.west)+(-0.12,0)$) {$C$};
\node[tinylabel, anchor=east] at ($(mat-4-1.west)+(-0.12,0)$) {$D$};
\foreach \c in {1,2,3,4} {
  \node[tinylabel, anchor=south] at ($(mat-1-\c.north)+(0,0.06)$) {$x_{\c}$};
}
\begin{scope}[on background layer]
  \fill[blue!10, rounded corners=1pt] ($(mat-1-1.north west)+(-0.04,0.04)$) rectangle ($(mat-2-4.south east)+(0.04,-0.04)$);
  \fill[blue!18, rounded corners=1pt] ($(mat-3-1.north west)+(-0.04,0.04)$) rectangle ($(mat-4-4.south east)+(0.04,-0.04)$);
\end{scope}

  \node[topanel] (p2) at (5.15,0) {};
\node[paneltitle] at ($(p2.north west)+(0.10,-0.12)$) {2. Build tree};
\node[
	anchor=north west,
	font=\scriptsize,
	align=left,
	text width=4.00cm,
] at ($(p2.north west)+(0.10,-0.38)$)
	{Hamming average-linkage hierarchy};

\node[treept] (r2) at ($(p2.center)+(0,1.35)$) {};
\node[treept] (u21) at ($(p2.center)+(-1.00,0.40)$) {};
\node[treept] (u22) at ($(p2.center)+(1.00,0.40)$) {};
\node[treept] (a2) at ($(p2.center)+(-1.55,-1.00)$) {};
\node[treept] (b2) at ($(p2.center)+(-0.45,-1.00)$) {};
\node[treept] (c2) at ($(p2.center)+(0.45,-1.00)$) {};
\node[treept] (d2) at ($(p2.center)+(1.55,-1.00)$) {};
\draw[treeline] (r2) -- (u21);
\draw[treeline] (r2) -- (u22);
\draw[treeline] (u21) -- (a2);
\draw[treeline] (u21) -- (b2);
\draw[treeline] (u22) -- (c2);
\draw[treeline] (u22) -- (d2);
\node[tinylabel, above=1pt of r2] {$R$};
\node[tinylabel, above=1pt of u21] {$U_1$};
\node[tinylabel, above=1pt of u22] {$U_2$};
\node[leafchip, anchor=north] at ($(a2)+(0,-0.24)$) {$A$\\$[1,1,0,1]$};
\node[leafchip, anchor=north] at ($(b2)+(0,-0.24)$) {$B$\\$[1,1,0,1]$};
\node[leafchip, anchor=north] at ($(c2)+(0,-0.24)$) {$C$\\$[0,0,1,0]$};
\node[leafchip, anchor=north] at ($(d2)+(0,-0.24)$) {$D$\\$[0,0,1,0]$};

  \node[widepanel] (p3) at (2.58,-5.32) {};
\node[paneltitle] at ($(p3.north west)+(0.10,-0.12)$) {3. Node distributions};
\node[
	anchor=north west,
	font=\scriptsize,
	align=left,
	text width=6.40cm,
] at ($(p3.north west)+(0.10,-0.38)$)
	{subtree-level distributions $\theta_u$ estimated from descendant leaves};
\node[
	anchor=north west,
	font=\scriptsize\itshape,
	align=left,
	text width=6.40cm,
] at ($(p3.north west)+(0.10,-0.78)$)
	{$\theta_u = \mathrm{mean}\{X_i : i \in L(u)\}$};

\node[treept] (r3) at ($(p3.center)+(0,1.42)$) {};
\node[treept] (u31) at ($(p3.center)+(-1.20,0.56)$) {};
\node[treept] (u32) at ($(p3.center)+(1.20,0.56)$) {};
\node[treept] (a3) at ($(p3.center)+(-1.92,-0.28)$) {};
\node[treept] (b3) at ($(p3.center)+(-0.48,-0.28)$) {};
\node[treept] (c3) at ($(p3.center)+(0.48,-0.28)$) {};
\node[treept] (d3) at ($(p3.center)+(1.92,-0.28)$) {};
\draw[treeline] (r3) -- (u31);
\draw[treeline] (r3) -- (u32);
\draw[treeline] (u31) -- (a3);
\draw[treeline] (u31) -- (b3);
\draw[treeline] (u32) -- (c3);
\draw[treeline] (u32) -- (d3);
\node[leafchip, anchor=north] at ($(a3)+(0,-0.24)$) {$A$\\$[1,1,0,1]$};
\node[leafchip, anchor=north] at ($(b3)+(0,-0.24)$) {$B$\\$[1,1,0,1]$};
\node[leafchip, anchor=north] at ($(c3)+(0,-0.24)$) {$C$\\$[0,0,1,0]$};
\node[leafchip, anchor=north] at ($(d3)+(0,-0.24)$) {$D$\\$[0,0,1,0]$};
\matrix[
	matrix of nodes,
	anchor=south,
	row sep=0pt,
	column sep=0.55cm,
	nodes={summarychip, anchor=center},
] (distrow) at ($(p3.south)+(0,0.82)$) {
	{$\theta_{U_1}=$\\$(1,1,0,1)$} & {$\theta_R=$\\$(.5,.5,.5,.5)$} & {$\theta_{U_2}=$\\$(0,0,1,0)$} \\
};
\draw[helperline] (distrow-1-1.north) -- ($(u31)+(0,-0.12)$);
\draw[helperline] (distrow-1-2.north) -- ($(r3)+(0,-0.14)$);
\draw[helperline] (distrow-1-3.north) -- ($(u32)+(0,-0.12)$);

\end{scope}

% Middle block: cards 4 and 5 are laid out side by side on one row.
\begin{scope}[shift={(0,-9.55)}]
  \node[panel, fill=orange!3] (p4) at (0,0) {};
  \node[paneltitle] at ($(p4.north west)+(0.10,-0.12)$) {4. Child-parent test};
  
\node[treept] (r4) at ($(p4.center)+(0,1.30)$) {};
\node[treept] (u41) at ($(p4.center)+(-1.05,0.52)$) {};
\node[treept] (u42) at ($(p4.center)+(1.05,0.52)$) {};
\node[treept] (a4) at ($(p4.center)+(-1.60,-0.30)$) {};
\node[treept] (b4) at ($(p4.center)+(-0.50,-0.30)$) {};
\node[treept] (c4) at ($(p4.center)+(0.50,-0.30)$) {};
\node[treept] (d4) at ($(p4.center)+(1.60,-0.30)$) {};
\draw[treeline] (r4) -- (u41);
\draw[treeline] (r4) -- (u42);
\draw[treeline] (u41) -- (a4);
\draw[treeline] (u41) -- (b4);
\draw[treeline] (u42) -- (c4);
\draw[treeline] (u42) -- (d4);
\draw[edgegateline] (r4) -- (u41);
\node[tinylabel, anchor=north] at ($(a4)+(0,-0.22)$) {$A$};
\node[tinylabel, anchor=north] at ($(b4)+(0,-0.22)$) {$B$};
\node[tinylabel, anchor=north] at ($(c4)+(0,-0.22)$) {$C$};
\node[tinylabel, anchor=north] at ($(d4)+(0,-0.22)$) {$D$};

\end{scope}

\begin{scope}[shift={(6.65,-9.55)}]
  \node[panel, fill=green!3] (p5) at (0,0) {};
  \node[paneltitle] at ($(p5.north west)+(0.10,-0.12)$) {5. Sibling test};
  \input{figures/pipeline_overview_modules/card5_sibling_test_body}
\end{scope}

% Bottom block: the final cut card is centered under the two test cards.
\begin{scope}[shift={(3.30,-14.85)}]
  \node[panelsmall] (p6) at (0,0) {};
\node[paneltitle] at ($(p6.north west)+(0.10,-0.12)$) {6. Final cluster cut};

\node[treept] (r6) at ($(p6.center)+(0,1.10)$) {};
\node[treept] (u61) at ($(p6.center)+(-0.94,0.18)$) {};
\node[treept] (u62) at ($(p6.center)+(0.94,0.18)$) {};
\node[treept] (a6) at ($(p6.center)+(-1.55,-0.95)$) {};
\node[treept] (b6) at ($(p6.center)+(-0.45,-0.95)$) {};
\node[treept] (c6) at ($(p6.center)+(0.45,-0.95)$) {};
\node[treept] (d6) at ($(p6.center)+(1.55,-0.95)$) {};
\draw[treeline] (r6) -- (u61);
\draw[treeline] (r6) -- (u62);
\draw[treeline] (u61) -- (a6);
\draw[treeline] (u61) -- (b6);
\draw[treeline] (u62) -- (c6);
\draw[treeline] (u62) -- (d6);
\draw[cutline] ($(r6)+(-0.22,0.42)$) -- ($(r6)+(-1.00,-0.18)$);
\draw[cutline] ($(r6)+(0.22,0.42)$) -- ($(r6)+(1.00,-0.18)$);
\begin{scope}[on background layer]
  \fill[blue!10, rounded corners=2pt] ($(u61)+(-0.68,-1.05)$) rectangle ($(u61)+(0.68,0.06)$);
  \fill[green!10, rounded corners=2pt] ($(u62)+(-0.68,-1.05)$) rectangle ($(u62)+(0.68,0.06)$);
\end{scope}
\node[leafchip, anchor=north] at ($(a6)+(0,-0.24)$) {$A$\\$[1,1,0,1]$};
\node[leafchip, anchor=north] at ($(b6)+(0,-0.24)$) {$B$\\$[1,1,0,1]$};
\node[leafchip, anchor=north] at ($(c6)+(0,-0.24)$) {$C$\\$[0,0,1,0]$};
\node[leafchip, anchor=north] at ($(d6)+(0,-0.24)$) {$D$\\$[0,0,1,0]$};
\node[tinylabel, anchor=north] at ($(p6.south)+(-0.95,0.22)$) {cluster 1};
\node[tinylabel, anchor=north] at ($(p6.south)+(0.95,0.22)$) {cluster 2};

\end{scope}

% Flow arrows between blocks.
\draw[steparrow] ($(p1.east)+(0.18,0)$) -- ($(p2.west)+(-0.18,0)$);
\draw[steparrow] ($(p2.south)+(-0.90,-0.05)$) -- ($(p3.north)+(1.85,0.10)$);
\draw[steparrow] ($(p3.south)+(-1.85,-0.03)$) -- ($(p4.north)+(0.40,0.10)$);
\draw[steparrow] ($(p3.south)+(1.85,-0.03)$) -- ($(p5.north)+(-0.40,0.10)$);
\draw[steparrow] ($(p4.south east)+(0.05,-0.04)$) -- ($(p6.north west)+(-0.12,0.06)$);
\draw[steparrow] ($(p5.south west)+(-0.05,-0.04)$) -- ($(p6.north east)+(0.12,0.06)$);

\end{tikzpicture}}
\caption{Overview of the KL-TE method. Samples are represented in a typed
sample-feature matrix, organized into a candidate hierarchy, summarized by
subtree-level distributions, evaluated with child-parent and sibling tests, and
finally converted into terminal clusters by a top-down traversal. The schematic
illustrates the inferential structure of the method rather than a literal
execution trace.}
\end{figure}
